\Needspace{18\baselineskip}
\section{The five search rules on underlying graph coordinates}
\label{sec:search}

The next operational dependency is Fig.~9(a--e). We now define these five rules
on actual ordered incidences of a fixed hypernet and prove determinism and
preservation under shallow hole filling. The token is represented by a mode
and a vertex of the underlying graph. This follows the paper's observation
that search leaves the underlying graph unchanged. The relationship to the
explicit focus-edge representation is still an adequacy obligation: explicit insertion and incidence recovery are now proved in
Section~\ref{sec:realization}, and Section~\ref{sec:erasure} supplies general
erasure and search on explicit states. Equivalence with the paper's contextual
graph rewrites remains open. The relation is
defined even on graphs with holes; its intended state interpretation restricts
the underlying graph to be focus-free, hole-free, and of type $*\to\epsilon$,
with a token on a term vertex. The stronger structural lemmas below do not
need these restrictions.

\subsection{Operation kinds and certified incidence views}
\begin{definition}[Operation kind]\label{term:operation-kind}
The two operation kinds are passive and active. A function from operations to
these kinds specifies the partition from \S5.1. Finishing a passive operation
produces backtrack mode; finishing an active operation produces rewrite mode.
Neither finishing mode is search, by distinction of constructors.
\end{definition}
\lean{Graph/Search}{10}{13}{Linear.OperationKind}
\lean{Graph/Search}{16}{18}{Linear.OperationKind.finish}
\lean{Graph/Search}{20}{21}{Linear.OperationKind.finish_ne_search}
The embedding of primitive labels into each depth alphabet is the universe lift
at zero and the left summand at successor depth. It is injective: project the
lift at depth zero, or use injectivity of the left injection at successor depth.
\lean{Graph/Search}{24}{26}{Linear.primitiveLabel}
\lean{Graph/Search}{28}{32}{Linear.primitiveLabel_injective}

\begin{definition}[Search operation view]\label{term:search-operation}
A search operation view records an actual edge $e$ with operation label $o$,
its source $s$, and lists $A,T$ of arguments and thunks. Its certificates assert
\[
 \mathsf{source}(e)=[s],\qquad
 \mathsf{target}(e)=A\mathbin{+\!+}T,\qquad
 |A|=\mathsf{arity}(o).
\]
Thus $A$ is exactly the ordered term-argument prefix. The ambient graph's
signature typing already specifies the types of the remaining thunk targets.
This view exposes certified data; it adds no machine-law assumption.
\end{definition}
\lean{Graph/Search}{36}{45}{Linear.SearchOperation}

\begin{lemma}[Existence and uniqueness of operation data]
Every shallow operation edge has such a view. Two views of the same edge have
the same operation, source, and argument list. Every argument is a target of
the edge, and the argument list has no repeated vertices.
\end{lemma}
\begin{proof}
Signature typing gives a singleton source list and a target list at least as
long as the arity. Choose its unique source; take the first arity-many targets
for $A$ and drop that prefix for $T$. The take/drop identity gives the target
certificate, and the length bound gives the arity certificate.
For uniqueness, the edge label determines the operation by injectivity of
primitive labels. Equality of singleton source lists determines the source.
The two argument lists have the same length; taking that many entries from the
equal target lists identifies them. Argument membership follows by inclusion
of the left append summand. No repetitions follows by restriction of the
whole target list's no-repetition certificate.
\end{proof}
\lean{Graph/Search}{52}{68}{Linear.SearchOperation.exists_view}
\lean{Graph/Search}{71}{82}{Linear.SearchOperation.data_unique}
\lean{Graph/Search}{84}{86}{Linear.SearchOperation.argument_mem}
\lean{Graph/Search}{88}{91}{Linear.SearchOperation.arguments_nodup}

\subsection{Local outcomes and the search relation}
\begin{definition}[Entry and return outcomes]\label{term:search-outcomes}
Entering an operation with argument list $A$ yields search at its first
argument if $A$ is nonempty. If $A$ is empty, it yields the operation's finishing
mode at its source. Returning from argument index $i$ yields search at argument
$i+1$ if that entry exists; otherwise it yields the finishing mode at the
source. Indices are zero-based. Return transitions additionally require that
argument $i$ exists, so the last case cannot be reached by returning from an
out-of-range index.
\end{definition}
\lean{Graph/Search}{94}{97}{Linear.SearchOperation.enter}
\lean{Graph/Search}{100}{103}{Linear.SearchOperation.resume}
These outcomes are independent of the chosen view of an edge. Substitute the
operation, source, and argument equalities from data uniqueness into their
definitions; the resulting expressions are identical.
\lean{Graph/Search}{105}{108}{Linear.SearchOperation.enter_unique}
\lean{Graph/Search}{110}{113}{Linear.SearchOperation.resume_unique}

\begin{definition}[Search position and step]\label{term:search-step}
A search position is a pair $(m,v)$ of mode and outer graph vertex. The search
relation has three constructors covering the five rules:
\begin{enumerate}
\item At a variable edge's singleton source $v$, search becomes rewrite at $v$
(Fig.~9a).
\item At an operation's source, search yields its entry outcome
(Fig.~9b for nonempty arguments and Fig.~9e for empty arguments).
\item Backtracking at argument $i$ yields its return outcome
(Fig.~9c when another argument exists and Fig.~9d when it does not).
\end{enumerate}
Every edge mentioned is shallow. Box contents and thunk targets are not searched
by these rules. No rule changes the underlying graph, and this relation contains
neither substitution nor extrinsic computation steps.
\end{definition}
\lean{Graph/Search}{151}{151}{Linear.SearchPosition}
\lean{Graph/Search}{155}{163}{Linear.SearchStep}
The four named operation rules specialize the two grouped constructors by
rewriting the argument list or the lookup of the next argument. Together with
the variable constructor, they expose precisely the five cases of Fig.~9(a--e).
\lean{Graph/Search}{169}{172}{Linear.SearchStep.first_argument}
\lean{Graph/Search}{175}{177}{Linear.SearchStep.nullary}
\lean{Graph/Search}{180}{183}{Linear.SearchStep.next_argument}
\lean{Graph/Search}{186}{189}{Linear.SearchStep.last_argument}

\subsection{Determinism and initial/final boundaries}
\begin{theorem}[Search determinism]
Two search steps with the same input coordinates have the same output
coordinates.
\end{theorem}
\begin{proof}
Case-split on their rule constructors. Search and backtrack inputs cannot agree.
For two variable steps the result is rewrite at the same vertex. For a variable
step and an operation-entry step, uniqueness of the edge having the given
source forces their edges equal, contradicting their distinct primitive labels.
For two operation-entry steps, source incidence uniqueness again identifies the
edges; entry-outcome uniqueness identifies the results.
For two return steps, both starting vertices occur in their operation target
lists. Target incidence uniqueness identifies the edges and hence their
argument lists. Since these lists have no repetitions, equal arguments force
equal indices. Return-outcome uniqueness then identifies the results.
The more general helper permits two separately named input coordinates with an
explicit equality; the deterministic theorem specializes that equality to
reflexivity. No determinism hypothesis is assumed.
\end{proof}
\lean{Graph/Search}{192}{232}{Linear.SearchStep.result_unique}
\lean{Graph/Search}{235}{237}{Linear.SearchStep.deterministic}

\begin{lemma}[Search boundary laws]
A rewrite-mode coordinate has no search step. A backtracking coordinate at an
input has no search step. No search step yields search mode at an input.
\end{lemma}
\begin{proof}
No constructor starts in rewrite mode. A step from backtrack must return from
an argument, which is a target of an edge. The graph's incidence conditions
forbid an input vertex from being a target.
For the last assertion, a variable step ends in rewrite, and finishing an
operation never produces search. Therefore a search-mode result must be an
argument selected on entry or on return. It is again an edge target and cannot
be an input. Apply this stronger non-input assertion to rule out any proposed
step into an initial coordinate.
\end{proof}
\lean{Graph/Search}{240}{242}{Linear.SearchStep.no_rewrite}
\lean{Graph/Search}{245}{249}{Linear.SearchStep.final_no_step}
\lean{Graph/Search}{252}{265}{Linear.SearchStep.search_target_not_input}
\lean{Graph/Search}{268}{270}{Linear.SearchStep.no_initial_target}
The finality assertion here concerns search steps only. It is not yet the
paper's assertion that final states have no transition of any kind; that also
requires the substitution and compute rules and the focus-edge bridge.

\subsection{Preservation under context filling}
\begin{lemma}[Transport of retained operation data]
Replacing an edge other than an operation edge retains a view of that operation.
Entry and return outcomes commute with the host port map.
\end{lemma}
\begin{proof}
Keep the operation label and use the surviving-edge injection. Map the source,
argument list, and thunk list through the host map. Mapping a singleton and an
append gives the two incidence certificates, and mapping preserves argument
length. Entry commutation follows by the empty/nonempty cases of the argument
list. Return commutation follows because lookup in a mapped list is the mapped
optional lookup; split according to whether the next argument exists. Modes
are unchanged by the port map.
\end{proof}
\lean{Graph/Search}{116}{127}{Linear.SearchOperation.plug}
\lean{Graph/Search}{129}{136}{Linear.SearchOperation.plug_enter}
\lean{Graph/Search}{138}{145}{Linear.SearchOperation.plug_resume}

\begin{theorem}[Existing context search steps survive filling]
Suppose $G$ has a shallow hole edge $h$ and a search step $p\to q$. For every
interface-compatible filling $H$, the graph $G[h:=H]$ has a search step from
the host image of $p$ to the host image of $q$.
\end{theorem}
\begin{proof}
In every rule, the interacting edge is a variable or operation and therefore
cannot equal the hole edge. In the variable case, retain that edge, its label,
and its mapped singleton source to apply the same rule. In the entry case,
use the retained operation view and entry commutation. In the return case,
lookup commutes with mapping, so the old argument witness supplies the mapped
argument witness; apply return commutation to identify the result.
\end{proof}
\lean{Graph/Search}{274}{308}{Linear.SearchStep.plug}
This theorem starts with a step of the context itself. Lemma~39(3) instead
starts with a step of a filled specimen and must establish that it comes from
the context. Extracting that step requires the entering/exiting analysis,
rootedness, output-closure, and the focus-edge bridge. The preservation theorem
proves one needed direction of locality, not the whole lemma.

\Needspace{24\baselineskip}
\subsection{Concrete checks of rule selection}
The fixture has one operation with two term arguments and one thunk. Its graph
is a corolla with targets $0,1,2$ in that order; the operation view certifies
that the argument list is $[0,1]$ and the thunk list is $[2]$. The incidence
certificates are checked by computation.
\lean{Graph/SearchExamples}{9}{9}{Linear.SearchExamples.ops}
\lean{Graph/SearchExamples}{11}{12}{Linear.SearchExamples.graph}
\lean{Graph/SearchExamples}{14}{23}{Linear.SearchExamples.operation}
Entry selects target $0$ and return from target $0$ selects target $1$. Return
from target $1$ finishes at the source in the declared mode, for either operation
kind. These are direct applications of the corresponding rule theorems.
\lean{Graph/SearchExamples}{26}{28}{Linear.SearchExamples.enter_first}
\lean{Graph/SearchExamples}{31}{33}{Linear.SearchExamples.return_next}
\lean{Graph/SearchExamples}{36}{38}{Linear.SearchExamples.finish_skips_thunk}
A hypothetical return step to search at thunk target $2$ would, by determinism,
have the same result as the finishing step. It would identify the finishing
mode with search, which is impossible. This checks that the full target list
has not accidentally been used as the argument list.
\lean{Graph/SearchExamples}{41}{45}{Linear.SearchExamples.cannot_search_thunk}
The existing nullary value corolla admits the empty argument view and finishes
by backtracking when declared passive. The unary variable corolla uses the
variable constructor and triggers rewrite without moving the coordinate.
\lean{Graph/SearchExamples}{48}{54}{Linear.SearchExamples.value_finishes}
\lean{Graph/SearchExamples}{57}{60}{Linear.SearchExamples.variable_triggers}
